\documentclass[../Main.tex]{subfiles}

\begin{document}
\section{Defining a Dynamical System}
\begin{definition}{Dynamical system}
    A \underline{dynamical system} is a mathematical entity described by a set of equations for a set of variables with respect to a time-like variable.
\end{definition}
In this course we will want to understand such systems. In particular, we will study:
\begin{itemize}
    \item long term behaviours of dynamical systems;
    \item periodic solutions;
    \item stability of solutions and chaos;
    \item how parameters of a system affect the solutions and behaviour.
\end{itemize}
\subsection{Continuous Time}
When our time-like variable is continuous, the dynamical system is described by a set of ODEs.
\begin{example}
    Consider a simple 1D system:
    \begin{equation*}
        \dot{x} = x(1-x^2)
    \end{equation*}
    This is a very easy example and in fact we find an analyical solution:
    \begin{equation*}
        x(t) = \pm \left(1-2e^2(t_0 - t)\right)^{-\frac12}
    \end{equation*}
    However, to practice the methods we would like to apply to more difficult systems, we will forget about this solution.

    We begin with analysis of fixed points. Solving the RHS gives stationary points at $x = 0, \pm 1$. Then we can also see that the RHS is positive on $(-\infty, -1)$ and $(0, 1)$, negative on the rest. This means that the fixed points at $-1$ and $1$ are stable (and attract solutions), whereas the fixed point at $0$ is unstable (and repels solutions). Therefore we can draw a diagram like in figure~\ref{figSimpleSystem}.
\end{example}
\begin{figure}
    \label{figSimpleSystem}
    \caption{Diagram of fixed points in a simple system}
\end{figure}
\begin{example}[Lotka-Volterra equations]
    Let $r$ be the number of rabbits in a population and $s$ be the number of sheep. Then for $a, b, c, \cdots, f > 0$, the Lotka-Volterra equations for these populations are:
    \begin{align*}
        \dot{r} &= r(a - br - cs) \\
        \dot{s} &= s(d - er - fs)
    \end{align*}
    The \underline{state space} of these equations is $(r, s) \in [0, \infty) \times [0, \infty)$.

    The constants $a$ and $d$ model unconstrained birth rate, and then $b, c, e, f$ model scarcity of food (where we assume that the rabbits and sheep both eat grass, which is constrained).
    Consider the example with $a = 3, d = 2, b=c=e=f=1$.
    \begin{align*}
        \dot{r} &= r(3 - r - s) \\
        \dot{s} &= s(2 - r - s)
    \end{align*}
    Then we can identify fixed points:
    \begin{equation*}
        (0, 0),~~(3, 0),~~(0, 2)
    \end{equation*}
    We also have fixed points for the individual equations when $r+s \in \{2, 3\}$.

    Now consider the signs of the derivatives. For large $r$ and $s$, both derivatives are negative. We can also see that $\dot{s} < 0$ when $r + s = 3$, and $\dot{r} > 0$ when $r + s = 2$. By then interpolating the rest, we get the diagram in figure~\ref{figLotkaVolterra}. We see that the behaviour as $t \to \infty$ is:
    \begin{equation*}
        (r, s) \to
        \begin{cases}
            (0, 0) & r = s = 0 \\
            (0, 2) & r = 0, s > 0 \\
            (3, 0) & \text{otherwise}.
        \end{cases}
    \end{equation*}
\end{example}
\begin{figure}
    \label{figLotkaVolterra}
    \caption{Diagram of solutions to Lotka-Volterra equations}
\end{figure}
\begin{example}[Predator-prey system]
    Consider a number of sheep $s$ and number of wolves $w$. These populations can be modelled by the system:
    \begin{align*}
        \dot{w} &= w(-a + bs) \\
        \dot{s} &= s(c - dw)
    \end{align*}
    for positive constants $a, b, c, d$.

    The fixed points are now $(0, 0)$ and $\left(\frac{a}{b}, \frac{c}{d}\right)$. By drawing our diagram, we see that we find periodic solutions around the non-zero fixed point.
\end{example}
\begin{example}
    Consider:
    \begin{align*}
        \dot{x} &= -y + \epsilon x(\mu - x^2 - y^2) \\
        \dot{y} &= x + \epsilon y (\mu - x^2 - y^2)
    \end{align*}
    We change coordinate system to plane polar coordinates to achieve:
    \begin{align*}
        \dot{r} &= \frac{x\dot{x} + y\dot{y}}{r} \\
        \dot{\theta} &= \frac{x\dot{y} - y\dot{x}}{r^2}
    \end{align*}
    then substituting this into the original equations gives:
    \begin{equation*}
        dot{r} = \epsilon r(\mu - r^2),\qquad\dot{\theta} = 1.
    \end{equation*}
    Now our system depends sensitively on $\epsilon$ and $\mu$. We can clearly see that $\epsilon = 0$ gives circular solutions, since $\dot{r} = 0$.

    In the case $\epsilon > 0, \mu \leq 0$ we have $\dot{r} < 0$ and the solution spirals in towards the origin. The case $\epsilon > 0, \mu \leq 0$ gives spirals outwards.

    The case $\mu > 0$ is more interesting. We now have a radial fixed point at $r = \sqrt{\mu}$.
    \begin{itemize}
        \item For $\epsilon > 0$, this fixed point is stable and the solution spirals towards it.
        \item For $\epsilon < 0$, this fixed point is unstable and the solution spirals away from it.
    \end{itemize}
\end{example}
\subsection{Discrete Time}
When our time-like variable is discrete, the dynamical system is described by a map. For example, the logistic map is:
\begin{equation*}
    x_{n+1} = \mu x_n(1-x_n),~~\mu \in [0, 1]
\end{equation*}
\end{document}